\documentclass[11pt]{article}
\usepackage{amsmath,amssymb}

\setlength{\parindent}{0pt}
\setlength{\parskip}{0.7em}
\setlength{\textwidth}{6.5in}
\setlength{\oddsidemargin}{0in}
\setlength{\evensidemargin}{0in}
\setlength{\topmargin}{-0.4in}
\setlength{\textheight}{9.2in}

\begin{document}

\begin{center}
{\Large SYSC 5108 Exam Study: Assignment 3, Question 1}\\
\vspace{0.3em}
{\large Entropy of a Bernoulli random variable and its maximum}\\
\vspace{0.3em}
Source question: \texttt{SYSC 5108 W26 Assignment 3.pdf}, Question 1
\end{center}

\section*{Question}

Maximum entropy is often used to select an a priori distribution when there is very little information about the random variable. Consider the Bernoulli distribution
\[
p(x)=
\begin{cases}
p, & x=1,\\
1-p, & x=0,\\
0, & \text{otherwise}.
\end{cases}
\]
Find its entropy $H(X)$ and prove analytically that $H(X)$ achieves its maximum when
\[
p=\frac{1}{2}.
\]

\section*{Step 1: Recall the entropy definition}

Chapter 3 slide 47 defines entropy as the expected information carried by a random variable:
\[
H(X)=-\sum_x p(x)\log p(x).
\]
Goodfellow, Bengio, and Courville define Shannon entropy the same way in Chapter 3:
\[
H(X)=\mathbb{E}_{X\sim P}[I(X)]
=-\mathbb{E}_{X\sim P}[\log P(X)].
\]
Unless otherwise specified, the course and textbook use natural logarithms, so the units are nats. If base-2 logarithms are used instead, the units are bits.

The slide convention
\[
0\log 0=0
\]
handles impossible outcomes.

\section*{Step 2: Apply the definition to a Bernoulli variable}

A Bernoulli random variable has two possible outcomes:
\[
P(X=1)=p,\qquad P(X=0)=1-p.
\]
So the entropy is
\[
H(X)
=-\left[P(X=1)\log P(X=1)+P(X=0)\log P(X=0)\right].
\]
Substitute the probabilities:
\[
\boxed{
H(X)=-p\log p-(1-p)\log(1-p)
}.
\]
This can also be written as
\[
H(X)=(p-1)\log(1-p)-p\log p.
\]

\section*{Step 3: Find the critical point}

Differentiate $H(p)$ with respect to $p$:
\[
H(p)=-p\log p-(1-p)\log(1-p).
\]
For the first term,
\[
\frac{d}{dp}\left[-p\log p\right]=-(\log p+1).
\]
For the second term,
\[
\frac{d}{dp}\left[-(1-p)\log(1-p)\right]
=\log(1-p)+1.
\]
Therefore
\[
\frac{dH}{dp}
=-(\log p+1)+\log(1-p)+1.
\]
Simplify:
\[
\frac{dH}{dp}
=\log(1-p)-\log p
=\log\left(\frac{1-p}{p}\right).
\]
Set the derivative equal to zero:
\[
\log\left(\frac{1-p}{p}\right)=0.
\]
Exponentiate both sides:
\[
\frac{1-p}{p}=1.
\]
So
\[
1-p=p.
\]
Hence
\[
\boxed{p=\frac{1}{2}.}
\]

\section*{Step 4: Prove the critical point is a maximum}

Compute the second derivative:
\[
\frac{dH}{dp}
=\log(1-p)-\log p.
\]
Differentiate again:
\[
\frac{d^2H}{dp^2}
=-\frac{1}{1-p}-\frac{1}{p}.
\]
Combine terms:
\[
\frac{d^2H}{dp^2}
=-\left(\frac{1}{1-p}+\frac{1}{p}\right).
\]
For $0<p<1$,
\[
\frac{1}{1-p}>0,\qquad \frac{1}{p}>0,
\]
so
\[
\frac{d^2H}{dp^2}<0.
\]
Therefore $H(p)$ is strictly concave on $(0,1)$, and the only critical point is the global maximum.

At $p=\frac12$:
\[
H\left(\frac12\right)
=-\frac12\log\frac12-\frac12\log\frac12.
\]
Since
\[
\log\frac12=-\log 2,
\]
we get
\[
H\left(\frac12\right)=\log 2.
\]
Thus the maximum entropy is
\[
\boxed{H_{\max}=\log 2\ \text{nats}.}
\]
If using base-2 logarithms, this is
\[
\boxed{H_{\max}=1\ \text{bit}.}
\]

\section*{Step 5: Check the boundary cases}

At $p=0$:
\[
H(0)=-0\log 0-(1)\log(1)=0.
\]
At $p=1$:
\[
H(1)=-(1)\log(1)-0\log 0=0.
\]
These are deterministic cases: if $p=0$, $X$ is always $0$; if $p=1$, $X$ is always $1$. There is no uncertainty, so entropy is zero.

At $p=\frac12$, both outcomes are equally likely:
\[
P(X=0)=P(X=1)=\frac12.
\]
This is the most uncertain Bernoulli distribution, so it has the maximum entropy.

\section*{Step 6: Exam Interpretation}

This result explains the maximum-entropy principle in the Bernoulli case. If there is no prior information favoring $X=0$ over $X=1$, choose the uniform Bernoulli distribution:
\[
P(X=0)=P(X=1)=\frac12.
\]
This is the least committed prior over the two outcomes because it has the largest entropy.

A concise exam answer could be:
\begin{quote}
For a Bernoulli random variable,
\[
H(X)=-\sum_xp(x)\log p(x)=-p\log p-(1-p)\log(1-p).
\]
Differentiate:
\[
\frac{dH}{dp}=\log(1-p)-\log p
=\log\left(\frac{1-p}{p}\right).
\]
Setting this to zero gives $(1-p)/p=1$, so $p=1/2$. The second derivative is
\[
\frac{d^2H}{dp^2}
=-\frac{1}{1-p}-\frac{1}{p}<0
\]
for $0<p<1$, so this critical point is a maximum. The maximum value is
\[
H(1/2)=\log 2
\]
nats, or $1$ bit if base-2 logarithms are used.
\end{quote}

\section*{Cross References Used}

\begin{itemize}
\item Assignment: \texttt{SYSC 5108 W26 Assignment 3.pdf}, Question 1.
\item Local submitted Assignment 3 PDF checked for comparison, Question 1 response.
\item Slides: Chapter 3 slide deck, slide 47 on entropy as average information and the convention $0\log0=0$.
\item Textbook: Goodfellow, Bengio, and Courville, \textit{Deep Learning}, Chapter 3, Section 3.9.1 on the Bernoulli distribution.
\item Textbook: Goodfellow, Bengio, and Courville, \textit{Deep Learning}, Chapter 3, Section 3.13 on information theory and Shannon entropy.
\end{itemize}

\end{document}
